% Ground-Truth-Anchored On-Policy Flow Matching -- draft v3 (received 2026-09-14).
%
% This is the CURRENT draft and supersedes the earlier rollout-based
% "Self-Forcing Flow Matching" write-up. It keeps the v2 idea (predict an
% endpoint, re-noise it, supervise with the REAL endpoint) but demotes it to the
% "cheap" variant of a wider compute--fidelity spectrum, and adds explicit
% K-step Euler rollouts plus a new supervision ablation (GT velocity vs GT
% endpoint) as first-class questions.
%
% Implementation status against this draft lives in
% docs/gt_anchored_on_policy_status.md.

\documentclass[11pt]{article}

\usepackage[margin=1in]{geometry}
\usepackage{amsmath,amssymb,mathtools}
\usepackage{booktabs}
\usepackage{hyperref}
\usepackage{enumitem}
\usepackage{xcolor}

\title{\textbf{Ground-Truth-Anchored On-Policy Flow Matching}\\
\large Learning from Model-Induced States at Controllable Rollout Cost}
\author{}
\date{}

\begin{document}
\maketitle

\begin{abstract}

Flow Matching (FM) trains a velocity field on states generated from a
prescribed coupling between data and noise. At inference time, however, the
learned vector field is evaluated along trajectories generated by the model
itself. Consequently, the states encountered during inference need not follow
the distribution of states used during FM training. Recent work has shown that
exposing generative models to model-induced states can alleviate related
train--inference mismatches, but doing so with faithful multi-step rollouts can
be computationally expensive.

We study a simple alternative: expose a Flow Matching model to states produced
by its own dynamics while retaining the original ground-truth data endpoint as
a corrective anchor. Our primary method performs an explicit one-step model
rollout from a conventional Flow Matching state. The resulting self-induced
state is detached and used as the input to a second model evaluation, while the
training target is constructed using the induced state and the original
ground-truth endpoint. This directly trains the model to make a corrective
prediction after its own dynamics have moved the state away from the prescribed
training path.

We further study increasingly faithful on-policy constructions using one, two,
and four model-induced steps. This provides a controlled compute--fidelity
tradeoff: deeper rollouts expose the model to states progressively farther
along its own dynamics, but require proportionally more model evaluations.
Finally, we study a cheap endpoint-based construction that predicts a model
endpoint and uses the model's corruption parameterization to construct a proxy
induced state without repeatedly executing the dynamics.

A central design question is the form of ground-truth supervision. We therefore
compare two alternatives: the original Flow Matching velocity target and a
ground-truth endpoint-directed corrective velocity computed from the
model-induced state. This distinguishes preserving the original transport
direction from explicitly correcting the model-induced state toward the true
data endpoint.

The resulting framework provides a controlled study of a broader question:
how much self-induced rollout is actually necessary to obtain the benefits of
on-policy Flow Matching adaptation?

\end{abstract}

\section{Introduction}

Flow Matching (FM) provides a simple framework for training continuous
generative vector fields. Given a coupling between a data sample $x_0$ and a
noise sample $\epsilon$, training states are generated according to a
prescribed path, commonly a straight interpolation. The velocity field is then
trained to match the velocity of this prescribed path.

At inference time, the situation is different. The model is no longer
evaluated on independently sampled states from the training coupling. Instead,
it generates a trajectory by repeatedly applying its own learned velocity
field. Thus, the distribution of states encountered during inference is
determined by the model itself.

This creates a potential train--inference state-distribution mismatch. A small
velocity error at one point in a trajectory can move the trajectory to a state
that is not well represented by the original Flow Matching training
distribution. The model must then make its next prediction from a state on
which it may have received relatively little direct supervision.

This issue is related to a broader class of train--test mismatch and exposure
bias problems. Reflow and related trajectory-rectification methods use
generated trajectories to modify the transport coupling. In autoregressive
video generation, Self-Forcing and the more recent Resampling Forcing expose
models to generated or degraded histories rather than exclusively to
ground-truth histories. MixFlow addresses a related training--testing
discrepancy for diffusion and Flow Matching models by modifying the training
state distribution. More recently, FlowCTS directly studies on-policy
trajectory supervision for Flow Matching models.

These works motivate a narrower question:

\begin{quote}
\emph{Can a Flow Matching model be improved simply by training it on states
created by its own dynamics, while retaining the original ground-truth
endpoint as a corrective anchor?}
\end{quote}

The distinction is important. We do not want to replace the ground-truth target
with the model's own potentially erroneous prediction. Instead, the model is
allowed to generate the \emph{input state}, while the original training
example continues to provide the \emph{target information}.

The central construction is therefore

\begin{equation}
\boxed{
\text{model-induced input state}
\quad+\quad
\text{ground-truth supervision}.
}
\end{equation}

Our primary method performs one explicit model dynamics step from a conventional
Flow Matching state. The resulting state is detached and passed through a
second model evaluation. The second evaluation is supervised using the
original ground-truth example.

We then investigate whether additional rollout depth is useful. Specifically,
we compare one-, two-, and four-step self-induced states. This allows us to
separate two questions that are often conflated:

\begin{enumerate}[leftmargin=*]
    \item whether exposure to model-induced states is useful at all; and
    \item how faithfully those states need to approximate the model's actual
    inference trajectory.
\end{enumerate}

Finally, we retain a cheaper endpoint-reconstruction construction from the
earlier version of the method. Rather than explicitly executing several model
dynamics steps, it predicts an endpoint and uses the model's corruption
parameterization to construct a proxy state. This creates a compute--fidelity
spectrum:

\begin{equation}
\boxed{
\text{cheap proxy}
\;\longleftrightarrow\;
\text{1-step}
\;\longleftrightarrow\;
\text{2-step}
\;\longleftrightarrow\;
\text{4-step}.
}
\end{equation}

The goal is not to assume that more rollout steps are always better. Instead,
we ask whether additional inference-state fidelity justifies its computational
cost.

\section{Problem Formulation}

Let a conditional flow model define the velocity field

\begin{equation}
\frac{d x_t}{dt}
=
v_\theta(x_t,t,c),
\end{equation}

where $c$ denotes conditioning information.

We use the convention that $t=1$ corresponds approximately to the noise
distribution and $t=0$ to the data distribution.

Standard Conditional Flow Matching constructs training states using a prescribed
coupling. For a straight path,

\begin{equation}
x_t = (1-t)x_0+t\epsilon,
\end{equation}

and trains the model with

\begin{equation}
\mathcal{L}_{\mathrm{FM}}
=
\mathbb{E}_{x_0,\epsilon,t}
\left[
\left\|
v_\theta(x_t,t,c)
-
u_t(x_0,\epsilon)
\right\|^2
\right].
\end{equation}

For the straight interpolation above, the path velocity is

\begin{equation}
u_t(x_0,\epsilon)=\epsilon-x_0.
\end{equation}

The optimal Flow Matching velocity can be written as

\begin{equation}
v_{\mathrm{FM}}^*(x,t,c)
=
\mathbb{E}
\left[
u_t(x_0,\epsilon)
\mid x_t=x,c
\right].
\end{equation}

During inference, however, the model generates its own ODE trajectory,

\begin{equation}
\frac{d x_t^\theta}{dt}
=
v_\theta(x_t^\theta,t,c).
\end{equation}

There is no guarantee that

\begin{equation}
p_{\mathrm{FM}}(x_t)
=
p_\theta^{\mathrm{ODE}}(x_t).
\end{equation}

If the learned velocity differs from the ideal velocity, the model can move
away from the prescribed training path. Subsequent velocity evaluations are
then performed at model-induced states.

Our objective is to directly expose the model to such states while maintaining
access to the original ground-truth training example.

\section{Ground-Truth-Anchored Self-Induced Training}

\subsection{Core Principle}

The central principle is

\begin{equation}
\boxed{
\text{model-induced input state}
\quad+\quad
\text{ground-truth supervision}.
}
\end{equation}

The distinction between the input and target is fundamental.

The input state is allowed to be produced by the imperfect model itself. The
supervision remains tied to the real data example from which the original FM
training state was constructed.

Given an induced state $x_s^\theta$, where $s>0$, one possible ground-truth
endpoint-directed velocity target is

\begin{equation}
\boxed{
u_{\mathrm{EP}}
(x_s^\theta,s)
=
\frac{x_0-x_s^\theta}{0-s}.
}
\end{equation}

The corresponding loss is

\begin{equation}
\mathcal{L}_{\mathrm{EP}}
=
\mathbb{E}
\left[
\left\|
v_\theta(x_s^\theta,s,c)
-
u_{\mathrm{EP}}(x_s^\theta,s)
\right\|^2
\right].
\end{equation}

This target should be interpreted as a corrective endpoint-directed signal.
It does not claim that the chord between $x_s^\theta$ and $x_0$ is the unique
true conditional Flow Matching velocity at the induced state.

\subsection{Alternative: Ground-Truth Flow Velocity}

There is another natural target.

For the original straight FM coupling,

\begin{equation}
u_{\mathrm{FM}}
=
\epsilon-x_0.
\end{equation}

One can therefore train the model at the induced state using

\begin{equation}
\mathcal{L}_{\mathrm{Vel}}
=
\mathbb{E}
\left[
\left\|
v_\theta(x_s^\theta,s,c)
-
(\epsilon-x_0)
\right\|^2
\right].
\end{equation}

This target has a different interpretation from the endpoint-directed target.

The ground-truth velocity asks:

\begin{quote}
\emph{What was the velocity of the original prescribed FM trajectory?}
\end{quote}

The endpoint-directed target asks:

\begin{quote}
\emph{Given the model's current induced state, what constant velocity would
move this state toward the correct data endpoint?}
\end{quote}

These targets are not generally equivalent because $x_s^\theta$ need not lie on
the original straight training path.

We therefore treat

\begin{equation}
\boxed{
\text{GT velocity vs. GT endpoint}
}
\end{equation}

as a central design ablation rather than assuming one target is universally
correct.

\subsection{Why Not Simply Reuse the Original FM Velocity?}

Suppose an ordinary FM state $x_t$ is moved using the model and produces a new
state $x_s^\theta$.

The original FM velocity

\begin{equation}
u_t(x_0,\epsilon)
\end{equation}

is associated with the prescribed training path and the original state $x_t$.
After the model moves the state, the resulting $x_s^\theta$ need not lie on
that same prescribed straight path.

Therefore, using the original velocity target at the induced state implicitly
treats the model-induced state as though it were still an ordinary FM state.

This may nevertheless be a useful form of supervision. The GT-velocity
experiment therefore tests an alternative supervision hypothesis, while the
endpoint target changes the target according to the new state:

\begin{equation}
x_s^\theta
\longrightarrow
x_0.
\end{equation}

\section{One-Step Self-Induced Training}

\subsection{Generating the Induced State}

We begin with an ordinary FM training state

\begin{equation}
x_t^{\mathrm{FM}}
=
(1-t)x_0+t\epsilon.
\end{equation}

The model first evaluates

\begin{equation}
v_1
=
v_\theta(x_t^{\mathrm{FM}},t,c).
\end{equation}

We then take one explicit Euler step toward the data endpoint. Let

\begin{equation}
s=t-\Delta t.
\end{equation}

The induced state is

\begin{equation}
\boxed{
x_s^\theta
=
x_t^{\mathrm{FM}}
-
\Delta t\,v_1.
}
\end{equation}

For the initial method, $x_s^\theta$ is detached from the computation graph.

The model is then evaluated again:

\begin{equation}
v_2
=
v_\theta(x_s^\theta,s,c).
\end{equation}

Using the endpoint-directed target gives

\begin{equation}
u_{\mathrm{EP}}
=
\frac{x_0-x_s^\theta}{0-s}.
\end{equation}

The one-step loss is therefore

\begin{equation}
\boxed{
\mathcal{L}_{1\mathrm{-step}}
=
\left\|
v_\theta(x_s^\theta,s,c)
-
\frac{x_0-x_s^\theta}{0-s}
\right\|^2.
}
\end{equation}

The computational pattern is

\begin{equation}
\boxed{
x_t^{\mathrm{FM}}
\xrightarrow{\;v_\theta\;}
x_s^\theta
\xrightarrow{\;v_\theta\;}
v_2
\xrightarrow{\;x_0\;}
\mathcal{L}.
}
\end{equation}

The first forward pass creates the model-induced state. The second forward
pass receives that state and is directly optimized.

\subsection{Why Detach the First Step?}

The initial method deliberately detaches the induced state:

\begin{equation}
x_s^\theta
=
\operatorname{stopgrad}
\left(
x_t^{\mathrm{FM}}
-
\Delta t\,v_\theta(x_t^{\mathrm{FM}},t,c)
\right).
\end{equation}

Thus gradients are applied through the second velocity evaluation but not
through the dynamics that generated the training state.

This has two advantages.

First, it isolates the scientific question: does training on states generated
by the current model improve the velocity field?

Second, it avoids requiring differentiation through the rollout dynamics.
Gradient-through-rollout training can be investigated as a separate extension,
but is not necessary for the central hypothesis.

\subsection{Why One Step Can Matter}

The displacement created by one step is

\begin{equation}
\delta x
=
-\Delta t\,v_\theta(x_t,t,c).
\end{equation}

If the model has velocity error

\begin{equation}
e_v
=
v_\theta-v^*,
\end{equation}

then the resulting state error relative to ideal dynamics is approximately

\begin{equation}
\delta x_{\mathrm{error}}
\approx
-\Delta t\,e_v.
\end{equation}

Even when this displacement is small, the velocity field is evaluated at a
different point. Locally,

\begin{equation}
v_\theta(x+\delta x,t)
\approx
v_\theta(x,t)
+
J_v(x,t)\delta x,
\end{equation}

where $J_v$ is the spatial Jacobian of the velocity field.

Thus a small state displacement can produce a meaningful change in the
velocity query, particularly where the velocity field varies rapidly.

Repeated inference steps can accumulate such local errors, motivating the
multi-step experiments below.

\section{Multi-Step Self-Induced Training}

The one-step method can be generalized by repeatedly applying the model's own
dynamics before computing the training loss.

Starting from

\begin{equation}
x_t^{(0)}
=
x_t^{\mathrm{FM}},
\end{equation}

define

\begin{equation}
x_{t-\Delta t}^{(1)}
=
x_t^{(0)}
-
\Delta t
v_\theta(x_t^{(0)},t,c),
\end{equation}

and recursively,

\begin{equation}
x_{t-k\Delta t}^{(k)}
=
x_{t-(k-1)\Delta t}^{(k-1)}
-
\Delta t
v_\theta
\left(
x_{t-(k-1)\Delta t}^{(k-1)},
t-(k-1)\Delta t,
c
\right).
\end{equation}

The rollout states are detached before the final supervised evaluation.

For a $K$-step construction, the final state is

\begin{equation}
x_s^{(K)},
\qquad
s=t-K\Delta t.
\end{equation}

The endpoint-directed target is

\begin{equation}
u_{\mathrm{EP}}^{(K)}
=
\frac{x_0-x_s^{(K)}}{0-s}.
\end{equation}

Therefore,

\begin{equation}
\boxed{
\mathcal{L}_{K\mathrm{-step}}
=
\left\|
v_\theta(x_s^{(K)},s,c)
-
\frac{x_0-x_s^{(K)}}{0-s}
\right\|^2.
}
\end{equation}

We focus on

\begin{equation}
K\in\{1,2,4\}.
\end{equation}

These settings provide a simple progression from a local self-induced
perturbation to states increasingly far along the model's own dynamics.

\subsection{The Compute--Fidelity Tradeoff}

A central motivation for studying multiple rollout depths is that on-policy
fidelity and computational cost are directly coupled.

Ignoring implementation details, a $K$-step training example requires
approximately $K$ model evaluations to generate the induced state plus one
additional evaluation for the supervised prediction.

Thus,

\begin{equation}
C_K \propto K+1.
\end{equation}

\begin{center}
\begin{tabular}{c c c}
\toprule
Method & Rollout steps & Approx. forwards/example \\
\midrule
Standard FM & 0 & 1 \\
One-step & 1 & 2 \\
Two-step & 2 & 3 \\
Four-step & 4 & 5 \\
\bottomrule
\end{tabular}
\end{center}

This creates an important practical question:

\begin{quote}
\emph{How much additional benefit is obtained from more faithful model-induced
states per additional model evaluation?}
\end{quote}

It is possible that deeper rollouts provide progressively better adaptation.
However, it is also possible that most of the useful correction occurs after
the first model-induced deviation.

\section{Cheap Endpoint-Reconstruction Variant}

The explicit one-step and multi-step methods require repeatedly executing the
model dynamics. We therefore retain a cheaper construction based on the
model's predicted endpoint.

\subsection{Endpoint Prediction}

Given a conventional FM state,

\begin{equation}
x_t^{\mathrm{FM}},
\end{equation}

the model produces an estimate of the data endpoint,

\begin{equation}
x_t^{\mathrm{FM}}
\xrightarrow{\;v_\theta\;}
\hat{x}_0^\theta.
\end{equation}

The predicted endpoint is detached.

Rather than executing a multi-step ODE rollout, we use the model's
corruption/noising parameterization to construct a proxy intermediate state:

\begin{equation}
\boxed{
\hat{x}_{t'}^\theta
=
\operatorname{Corrupt}
\left(
\hat{x}_0^\theta,\epsilon',t'
\right).
}
\end{equation}

For a straight interpolation parameterization this can take the form

\begin{equation}
\hat{x}_{t'}^\theta
=
(1-t')\hat{x}_0^\theta
+
t'\epsilon'.
\end{equation}

The crucial point is that $\hat{x}_{t'}^\theta$ is treated as a model-dependent
training input, while the original endpoint remains

\begin{equation}
x_0.
\end{equation}

We therefore use

\begin{equation}
u_{\mathrm{cheap}}
=
\frac{x_0-\hat{x}_{t'}^\theta}{0-t'}
\end{equation}

and train with

\begin{equation}
\boxed{
\mathcal{L}_{\mathrm{cheap}}
=
\left\|
v_\theta(\hat{x}_{t'}^\theta,t',c)
-
\frac{x_0-\hat{x}_{t'}^\theta}{0-t'}
\right\|^2.
}
\end{equation}

\subsection{Interpretation}

The cheap construction is not identical to true model dynamics.

The explicit rollout method creates

\begin{equation}
x_t^{\mathrm{FM}}
\rightarrow
x_s^\theta
\end{equation}

through the learned vector field.

The cheap method instead creates

\begin{equation}
x_t^{\mathrm{FM}}
\rightarrow
\hat{x}_0^\theta
\rightarrow
\hat{x}_{t'}^\theta
\end{equation}

through endpoint prediction followed by the corruption parameterization.

Therefore, the cheap state should be viewed as a \emph{proxy for a
model-dependent inference state}, rather than as an exact ODE state.

This distinction gives rise to an important empirical question:

\begin{quote}
\emph{Does faithful multi-step self-induction provide substantially more useful
training states than a much cheaper endpoint-based proxy?}
\end{quote}

\section{Combining with Standard Flow Matching}

A conservative objective can retain the original FM loss during post-training:

\begin{equation}
\boxed{
\mathcal{L}
=
\lambda\mathcal{L}_{\mathrm{FM}}
+
(1-\lambda)\mathcal{L}_{\mathrm{on}}.
}
\end{equation}

Here $\mathcal{L}_{\mathrm{on}}$ can correspond to

\begin{equation}
\mathcal{L}_{1\mathrm{-step}},
\quad
\mathcal{L}_{2\mathrm{-step}},
\quad
\mathcal{L}_{4\mathrm{-step}},
\end{equation}

or

\begin{equation}
\mathcal{L}_{\mathrm{cheap}}.
\end{equation}

The value of retaining the original FM objective should be evaluated
experimentally rather than assumed to be necessary.

\section{Relation to Inference-Time Rollouts}

The idealized version of the hypothesis would directly sample states from a
complete inference trajectory,

\begin{equation}
x_t^\theta
\rightarrow
x_{t-\Delta t}^\theta
\rightarrow
x_{t-2\Delta t}^\theta
\rightarrow
\cdots
\rightarrow
x_0^\theta.
\end{equation}

The one-step method approximates this process locally:

\begin{equation}
x_t^{\mathrm{FM}}
\rightarrow
x_s^\theta.
\end{equation}

The two- and four-step methods progressively extend this approximation.

The cheap endpoint method makes a different approximation:

\begin{equation}
x_t^{\mathrm{FM}}
\rightarrow
\hat{x}_0^\theta
\rightarrow
\hat{x}_{t'}^\theta.
\end{equation}

Consequently, the methods form a controlled spectrum:

\begin{equation}
\boxed{
\text{cheap proxy}
\quad
\longleftrightarrow
\quad
\text{1-step}
\quad
\longleftrightarrow
\quad
\text{2-step}
\quad
\longleftrightarrow
\quad
\text{4-step}.
}
\end{equation}

The computational cost increases from left to right, while the fidelity to
the actual model dynamics is expected to increase in the same direction.

\section{Related Work}

\subsection{Flow Matching}

Flow Matching trains vector fields using prescribed probability paths and
provides the basic training framework considered here
\cite{lipman2023fm}.

\subsection{Rectified Flow and Reflow}

Rectified Flow and subsequent Reflow-style methods use generated trajectories
to construct new transport couplings and progressively improve the geometry of
the learned transport path \cite{liu2023rectified}.

Our method shares the use of model-generated states but has a different
objective. Rather than replacing the original training coupling with
generated endpoint pairs, we retain the original ground-truth example and use
it as an anchor while modifying the state presented to the model.

\subsection{Self-Forcing}

Self-Forcing addresses exposure bias in autoregressive video diffusion by
training models on their own generated contexts rather than relying exclusively
on ground-truth histories \cite{huang2025selfforcing}.

The common principle is exposure to inference-like states. Our setting is
different: the model is a continuous flow field, and the self-generated
quantity is an intermediate state along the model's generative dynamics rather
than an autoregressive history.

\subsection{Resampling Forcing}

Resampling Forcing is particularly relevant to the conceptual motivation.
It introduces a self-resampling procedure that simulates inference-time model
errors on history frames during training, allowing an autoregressive video
diffusion model to train on degraded/model-dependent histories while retaining
ground-truth targets \cite{guo2025resampling}.

The analogy is:

\begin{equation}
\text{model-dependent input}
+
\text{ground-truth target}.
\end{equation}

Our method transfers this broad principle to continuous Flow Matching
dynamics. Instead of generating an imperfect autoregressive history, the model
generates an intermediate state using its own velocity field. The key
additional question is then what ground-truth velocity signal should be used
at that induced state.

\subsection{MixFlow}

MixFlow studies the training--testing discrepancy in diffusion and Flow
Matching models and proposes to modify the training state distribution using
slowed interpolations that better reflect generated noisy states
\cite{li2025mixflow}.

MixFlow therefore provides a direct precedent for the importance of the
training-state distribution in Flow Matching. Our approach differs in that the
training state is explicitly generated by the current model's own dynamics,
and the original paired ground-truth endpoint is retained as corrective
supervision.

\subsection{On-Policy Distillation for Flow Models}

Flow-OPD and related on-policy distillation methods apply teacher-based
supervision to model-generated states or trajectories. These methods establish
that on-policy states can be useful for adapting flow models.

Our method is complementary in that the corrective signal is obtained directly
from the paired training example rather than requiring a separate teacher
trajectory.

\subsection{Self- and Teacher-Free On-Policy Flow Training}

Self-OPD and Any-OPD investigate teacher-free or heterogeneous on-policy
distillation settings for Flow Matching models. These approaches further
support the relevance of model-generated states for post-training.

Our focus is narrower: we ask whether a particularly simple ground-truth
anchor---the original data endpoint---can be used to correct model-induced
states without introducing a separate distillation teacher.

\subsection{FlowCTS}

FlowCTS is the most directly related recent work to the explicit multi-step
component of our method. It performs on-policy continuous trajectory
supervision for Flow Models by initializing subsequent student and reference
trajectories from a student-visited state. It also explicitly studies the
number of supervision steps and reports a tradeoff between richer trajectory
information and optimization difficulty \cite{ye2026flowcts}.

Our work differs in the source and interpretation of supervision. Rather than
matching a reference trajectory initialized from a student state, we retain
the original paired ground-truth endpoint and construct a corrective velocity
target from the induced state. Our main design question is consequently
different: how much explicit self-induced rollout is required, and can a cheap
endpoint-based proxy recover the benefit of deeper rollout?

\subsection{Partial-Path Rectification}

Partial-Path Rectification addresses accumulated sampling errors by
constructing ground-truth one-to-one image pairs for portions of the transport
path and enforcing consistency between subpaths \cite{zhang2026partial}.

The work is relevant because it demonstrates that ground-truth-anchored
correction can be useful for reducing errors accumulated along generative
trajectories. Our method instead focuses on adapting the model to states
generated by its own dynamics during training.

\subsection{Self-Distillation and Flow Post-Training}

Recent self-distillation and on-policy optimization approaches further
demonstrate the usefulness of training on model-generated states. The present
work is deliberately simpler: it asks whether the combination of

\begin{equation}
\text{self-induced state}
+
\text{ground-truth anchor}
\end{equation}

is itself a useful training primitive.

\section{Experimental Questions}

The paper is organized around six questions.

\subsection{Question 1: Do Model-Induced States Help?}

Compare standard FM post-training against ground-truth-anchored
self-induced training:

\begin{equation}
\mathrm{FM}
\quad\text{vs.}\quad
\mathrm{GT\text{-}anchored\ training}.
\end{equation}

The evaluation should measure generation quality and inference-time sampling
behavior.

\subsection{Question 2: Is Explicit Self-Induction Better Than a Proxy?}

Compare

\begin{equation}
\mathrm{Cheap}
\quad\text{vs.}\quad
\mathrm{1\text{-}step}.
\end{equation}

This distinguishes actual exposure to model dynamics from a cheaper
model-dependent state construction.

\subsection{Question 3: Does More Rollout Depth Help?}

Evaluate

\begin{equation}
K\in\{1,2,4\}.
\end{equation}

The important measurement is not only final quality but improvement per unit
compute.

\subsection{Question 4: What Kind of Ground-Truth Target Is Best?}

This is a central design ablation.

At the same induced state $x_s^\theta$, compare:

\paragraph{GT velocity:}

\begin{equation}
u_{\mathrm{FM}}
=
\epsilon-x_0,
\end{equation}

against

\paragraph{GT endpoint-directed velocity:}

\begin{equation}
u_{\mathrm{EP}}
=
\frac{x_0-x_s^\theta}{0-s}.
\end{equation}

The comparison asks whether adaptation benefits more from preserving the
original FM transport direction or from explicitly correcting the model's
current state toward the ground-truth endpoint.

\subsection{Question 5: Is Ground-Truth Anchoring Actually Necessary?}

Replace the ground-truth endpoint $x_0$ with the model's own endpoint
prediction $\hat{x}_0^\theta$.

This produces a self-consistency baseline:

\begin{equation}
x_s^\theta
\rightarrow
\hat{x}_0^\theta.
\end{equation}

The critical comparison is

\begin{equation}
\boxed{
\text{self-induced state + self endpoint}
\quad\text{vs.}\quad
\text{self-induced state + ground-truth endpoint}.
}
\end{equation}

If ground-truth anchoring produces a significant improvement, this supports
the interpretation that the method provides corrective supervision rather than
merely encouraging the model to reproduce its own behavior.

\subsection{Question 6: Is the Effect Specific to Model-Induced States?}

Perturb FM states using a generic displacement whose magnitude is matched to
the displacement produced by a model step.

If

\begin{equation}
\|\delta_{\mathrm{generic}}\|
\approx
\|x_s^\theta-x_t^{\mathrm{FM}}\|,
\end{equation}

then compare

\begin{equation}
x_s^\theta
\end{equation}

against

\begin{equation}
x_t^{\mathrm{FM}}+\delta_{\mathrm{generic}}.
\end{equation}

If self-induced states outperform generic perturbations, this provides
evidence that the benefit comes specifically from adapting to model-induced
inference states rather than simply performing state-space augmentation.

\section{Experimental Design}

\subsection{Baseline}

Begin with a pretrained Flow Matching model and perform post-training under
controlled compute budgets.

The standard FM baseline uses

\begin{equation}
\mathcal{L}_{\mathrm{FM}}
\end{equation}

on the original training-state distribution.

\subsection{One-Step Method}

For every training example:

\begin{enumerate}[leftmargin=*]
    \item Sample $(x_0,\epsilon,t)$.
    \item Construct $x_t^{\mathrm{FM}}$.
    \item Evaluate $v_\theta(x_t^{\mathrm{FM}},t)$.
    \item Take one Euler step to obtain $x_s^\theta$.
    \item Detach $x_s^\theta$.
    \item Evaluate the model at $(x_s^\theta,s)$.
    \item Construct the chosen ground-truth target.
    \item Backpropagate through the final evaluation.
\end{enumerate}

\subsection{Multi-Step Methods}

Repeat the state-generation process for

\begin{equation}
K=2
\quad\text{and}\quad
K=4.
\end{equation}

The final state is then used for the ground-truth-anchored loss.

\subsection{Cheap Method}

For the cheap variant:

\begin{enumerate}[leftmargin=*]
    \item Sample a conventional FM state.
    \item Predict $\hat{x}_0^\theta$.
    \item Detach the prediction.
    \item Construct $\hat{x}_{t'}^\theta$ using the corruption parameterization.
    \item Evaluate the model on $\hat{x}_{t'}^\theta$.
    \item Supervise using the original $x_0$.
\end{enumerate}

This avoids repeated explicit dynamics steps.

\subsection{Compute-Matched Evaluation}

All methods should be evaluated under both:

\begin{enumerate}[leftmargin=*]
    \item equal optimization steps; and
    \item approximately equal model-evaluation/FLOP budgets.
\end{enumerate}

The second comparison is particularly important because a four-step rollout
naturally performs substantially more forward evaluations per training example
than the cheap variant.

Results should report:

\begin{itemize}[leftmargin=*]
    \item generation quality;
    \item inference speed;
    \item training FLOPs or model evaluations;
    \item memory usage;
    \item quality as a function of training compute.
\end{itemize}

\section{Important Ablations}

\subsection{Rollout Depth}

\begin{equation}
K\in\{0,1,2,4\}.
\end{equation}

This tests whether deeper self-induced states progressively improve the model.

\subsection{Ground-Truth Velocity Versus Ground-Truth Endpoint}

At the same induced state, compare

\begin{equation}
\boxed{
u_{\mathrm{FM}}=\epsilon-x_0
}
\end{equation}

against

\begin{equation}
\boxed{
u_{\mathrm{EP}}
=
\frac{x_0-x_s^\theta}{0-s}.
}
\end{equation}

This should be one of the primary ablations because it determines what kind of
ground-truth information is most useful once the model has moved off the
prescribed training path.

\subsection{Ground-Truth Endpoint Versus Self-Predicted Endpoint}

Compare

\begin{equation}
x_0
\quad\text{vs.}\quad
\hat{x}_0^\theta.
\end{equation}

This tests whether the real-data anchor is responsible for the improvement.

\subsection{Self-Induced Versus Generic Perturbations}

Match the magnitude of a generic perturbation to the displacement created by
the model:

\begin{equation}
\|\delta_{\mathrm{generic}}\|
\approx
\|\delta_{\mathrm{self}}\|.
\end{equation}

This tests whether model-specific state exposure matters beyond generic
augmentation.

\subsection{Step Size}

Evaluate several $\Delta t$ values to determine whether the one-step benefit
depends strongly on the magnitude of the induced state displacement.

\subsection{Detach Versus Gradient Through Rollout}

The primary method uses

\begin{equation}
\operatorname{stopgrad}(x_s^\theta).
\end{equation}

An optional ablation can compare this against allowing gradients through the
state-generation step.

This should remain secondary because the detached formulation provides the
cleanest interpretation of the central hypothesis.

\subsection{FM Loss Weight}

Evaluate different values of $\lambda$ in

\begin{equation}
\mathcal{L}
=
\lambda\mathcal{L}_{\mathrm{FM}}
+
(1-\lambda)\mathcal{L}_{\mathrm{on}}.
\end{equation}

This determines whether retaining the original FM objective remains important
during post-training.

\section{Expected Results and Hypotheses}

We do not assume that deeper rollouts are universally superior. Instead, we
consider several possible outcomes.

\subsection{Outcome A: Progressive Improvement}

If

\begin{equation}
\mathrm{4\text{-}step}
>
\mathrm{2\text{-}step}
>
\mathrm{1\text{-}step}
>
\mathrm{FM},
\end{equation}

then increasingly faithful exposure to model-induced states appears beneficial.
The remaining question is whether the improvement justifies the additional
compute.

\subsection{Outcome B: Early Saturation}

If

\begin{equation}
\mathrm{1\text{-}step}
\approx
\mathrm{2\text{-}step}
\approx
\mathrm{4\text{-}step},
\end{equation}

then most of the useful correction may occur immediately after the model first
deviates from the prescribed FM state.

This would make one-step training particularly attractive because it provides
model-induced exposure at substantially lower cost than longer rollouts.

\subsection{Outcome C: Cheap Proxy Recovers Multi-Step Performance}

The most practically interesting result would be

\begin{equation}
\mathrm{Cheap}
\approx
\mathrm{2\text{-}step}
\approx
\mathrm{4\text{-}step},
\end{equation}

while requiring substantially fewer model evaluations.

This would indicate that explicit long rollouts are not necessary to capture
much of the useful adaptation.

Importantly, this would not imply that the cheap state is literally equivalent
to an inference state. Rather, it would demonstrate that a simpler
model-dependent proxy can provide a favorable approximation for training.

\subsection{Outcome D: Ground-Truth Endpoint Outperforms GT Velocity}

A particularly informative result would be

\begin{equation}
\mathcal{L}_{\mathrm{EP}}
<
\mathcal{L}_{\mathrm{Vel}}
\end{equation}

in terms of downstream generation quality.

This would suggest that once the model moves away from the prescribed FM path,
the most useful ground-truth information is not necessarily the velocity of the
original coupling, but rather information that explicitly accounts for the
model's current state and points it toward the true endpoint.

Conversely, if GT velocity performs better, this would suggest that preserving
the original FM vector field is more important than constructing a new
state-dependent corrective direction.

Either result would provide useful information about the mechanism.

\section{Optional Extensions}

The main paper should initially focus on the one-step, multi-step, and cheap
variants described above.

Several extensions can be explored if the basic effect is established.

\subsection{Longer Rollouts}

The rollout depth can be extended beyond four steps:

\begin{equation}
K\in\{1,2,4,8,\ldots\}.
\end{equation}

The purpose would be to determine whether performance eventually saturates or
degrades as the training state moves farther from the original FM coupling.

\subsection{Hierarchical Self-Induced Training}

A possible extension is to recursively generate training states from previously
induced states, producing a hierarchy of increasingly model-dependent states.

This is intentionally not part of the primary method because it introduces
additional complexity and makes it harder to isolate the contribution of the
basic one-step construction.

\subsection{Gradient-Through-Rollout Training}

Another extension is to remove the stop-gradient operation and optimize through
the state-generation dynamics.

This would transform the method from detached state adaptation toward a more
direct optimization of multi-step trajectory behavior. Again, this should be
treated as a later experiment rather than as part of the initial formulation.

\section{Discussion}

The proposed framework separates two questions that are often conflated.

The first is whether training on model-induced states is useful.

The second is how accurately those states need to reproduce the actual
inference distribution.

The one-step and multi-step experiments address the first question while
progressively testing the second.

A one-step state is not expected to reproduce a complete inference trajectory.
Instead, it represents a local intervention:

\begin{equation}
\text{prescribed training state}
\rightarrow
\text{state after model-induced deviation}.
\end{equation}

This may already be sufficient to teach the model how to behave after its own
predictions perturb the state.

The cheap variant tests an even broader hypothesis: perhaps exact trajectory
reproduction is unnecessary, provided that the training procedure generates
states that reflect the kinds of errors produced by the model.

The GT velocity versus GT endpoint experiment further separates the
\emph{state-generation} question from the \emph{supervision} question.

Thus the method can be viewed as the combination of two design axes:

\begin{equation}
\boxed{
\text{How is the training state generated?}
}
\end{equation}

and

\begin{equation}
\boxed{
\text{What ground-truth signal is used there?}
}
\end{equation}

The first axis ranges from cheap proxies to increasingly long explicit
rollouts. The second ranges from preserving the original FM velocity to
constructing a state-dependent endpoint correction.

This gives a compact experimental matrix:

\begin{center}
\begin{tabular}{l c c}
\toprule
State construction & GT velocity & GT endpoint \\
\midrule
FM state & baseline & baseline/extension \\
Cheap proxy & ablation & primary cheap method \\
1-step self-induced & primary ablation & primary method \\
2-step self-induced & ablation & primary method \\
4-step self-induced & ablation & primary method \\
\bottomrule
\end{tabular}
\end{center}

The goal is not simply to make training more on-policy at any cost. The goal is
to determine the minimum amount of self-induced computation and the appropriate
form of ground-truth supervision required to obtain useful adaptation.

\section{Limitations}

Several limitations should be acknowledged.

First, a one-step Euler update is only a local approximation to the model's
inference trajectory.

Second, the ground-truth endpoint-directed velocity is not necessarily equal
to the exact conditional Flow Matching velocity at an arbitrary model-induced
state. It is therefore best interpreted as corrective supervision.

Third, the original FM velocity is also not guaranteed to be the correct
velocity at an induced state. The GT-velocity experiment therefore tests an
alternative supervision hypothesis rather than providing an exact oracle.

Fourth, sufficiently large or repeated self-induced deviations may leave the
meaningful data region and produce harmful training states.

Fifth, the cheap endpoint-reconstruction state is not an exact on-policy ODE
state. Its usefulness must therefore be demonstrated empirically.

Finally, the existence of a train--inference mismatch does not by itself imply
that correcting this mismatch is the dominant limitation of Flow Matching.
The proposed experiments are intended to test whether explicitly addressing
this mismatch improves generation.

\section{Conclusion}

We propose Ground-Truth-Anchored On-Policy Flow Matching, a simple framework
for adapting Flow Matching models to states generated by their own dynamics
while retaining real ground-truth supervision.

The central construction is

\begin{equation}
\boxed{
\text{self-induced input state}
+
\text{ground-truth supervision}.
}
\end{equation}

Our primary method performs one explicit model step from an ordinary FM state,
then trains the model on the resulting state using ground-truth information
associated with the original data example.

We further study two-, four-, and potentially longer self-induced rollouts to
measure the tradeoff between inference-state fidelity and computational cost.
This produces a controlled spectrum of increasingly expensive on-policy
training.

Finally, we study a cheap endpoint-based approximation that predicts the data
endpoint once and constructs a proxy model-dependent state without repeatedly
executing expensive model dynamics.

A central new design question is whether the correct ground-truth supervision
at an induced state should be the original FM velocity,

\begin{equation}
\epsilon-x_0,
\end{equation}

or a state-dependent corrective velocity,

\begin{equation}
\frac{x_0-x_s^\theta}{0-s}.
\end{equation}

The central empirical question is therefore:

\begin{quote}
\emph{How much self-induced computation, and what form of ground-truth
supervision, are actually necessary to obtain the benefits of on-policy
Flow Matching adaptation?}
\end{quote}

If performance improves monotonically with rollout depth, the method provides a
simple way to trade compute for increasingly faithful inference-state
exposure. If performance saturates after one step, local self-induced training
may provide a particularly efficient alternative to full rollouts. If the
cheap endpoint-based construction approaches the performance of multi-step
training, it would show that useful on-policy adaptation can be achieved
without repeatedly executing expensive model dynamics.

\begin{thebibliography}{20}

\bibitem{lipman2023fm}
Y. Lipman, R. T. Q. Chen, H. Ben-Hamu, M. Nickel, and M. Le.
\newblock Flow Matching for Generative Modeling.
\newblock ICLR, 2023.

\bibitem{liu2023rectified}
X. Liu, C. Gong, and Q. Liu.
\newblock Flow Straight and Fast: Learning to Generate and Transfer Data with Rectified Flow.
\newblock ICLR, 2023.

\bibitem{huang2025selfforcing}
X. Huang et al.
\newblock Self Forcing: Bridging the Train-Test Gap in Autoregressive Video Diffusion.
\newblock NeurIPS, 2025.
\newblock arXiv:2506.08009.

\bibitem{guo2025resampling}
Y. Guo et al.
\newblock End-to-End Training for Autoregressive Video Diffusion via Self-Resampling.
\newblock ECCV, 2026.
\newblock arXiv:2512.15702.

\bibitem{li2025mixflow}
H. Li, J. Lyu, F.-Y. Wang, K. Cheng, S. Zhu, and J. Wang.
\newblock MixFlow Training: Alleviating Exposure Bias with Slowed Interpolation Mixture.
\newblock CVPR, 2026.
\newblock arXiv:2512.19311.

\bibitem{fang2026flowopd}
Z. Fang et al.
\newblock Flow-OPD: On-Policy Distillation for Flow Matching Models.
\newblock arXiv preprint arXiv:2605.08063, 2026.

\bibitem{zhang2026selfopd}
S. Zhang et al.
\newblock Self-OPD: On-Policy Distillation for Flow Matching Models without Teacher.
\newblock arXiv preprint arXiv:2608.26872, 2026.

\bibitem{fu2026anyopd}
S. Fu et al.
\newblock Any-OPD: Heterogeneous On-Policy Distillation for Flow-Matching Models via Representation-Space Bridging.
\newblock arXiv preprint arXiv:2608.03316, 2026.

\bibitem{ye2026flowcts}
K. Ye et al.
\newblock FlowCTS: On-policy Continuous Trajectory Supervision of Flow Models.
\newblock arXiv preprint arXiv:2607.24522, 2026.

\bibitem{zhang2026partial}
J. Zhang, D. Liu, J. H. Ko, E. Park, S. Zhang, and C. Xu.
\newblock Partial-Path Rectification in Diffusion Sampling.
\newblock IEEE Transactions on Multimedia, 2026.

\bibitem{roy2026selfdistill}
S. Roy, D. Mukherjee, and P. Patil.
\newblock Optimal Self-Distillation for Rectified Flow via Linear Probing.
\newblock arXiv preprint arXiv:2607.14947, 2026.

\bibitem{porcher2026flowwm}
F. Porcher, N. Carion, K. Alahari, and S. Chen.
\newblock Flow Matching in Feature Space for Stochastic World Modeling.
\newblock arXiv preprint arXiv:2606.29059, 2026.

\end{thebibliography}

\end{document}
